\begin{lemma}[Artin-Rees]
\label{lemma-Artin-Rees}
Suppose that $R$ is Noetherian. Let $I \subset R$ be an ideal.
Let $N \subset M$ be finite $R$-modules.
There exists a constant $c > 0$ such that
$I^n M \cap N  =  I^{n-c}(I^cM \cap N)$ for all $n \geq c$.
\end{lemma}

\begin{proof}
Keep the graded ring
$$
S=\bigoplus_{n\geq0}I^n,\qquad I^0=R,
$$
with its actual degree inclusions $\iota_n:I^n\to S$.
Thus $\iota_n(a)\iota_m(b)=\iota_{n+m}(ab)$.
Since $R$ is Noetherian, write $I=(f_1,\ldots,f_t)$.
The $R$-algebra map $R[X_1,\ldots,X_t]\to S$ sends $r\in R$
to $\iota_0(r)$ and sends a monomial with its coefficient to
$$
rX_1^{e_1}\cdots X_t^{e_t}
\longmapsto
\iota_{e_1+\cdots+e_t}
 (rf_1^{e_1}\cdots f_t^{e_t}).
$$
Every element of $I^n$ is a finite sum of these products of total
degree $n$, so the map is surjective. The polynomial ring, and
therefore $S$, is Noetherian by Lemma
\ref{lemma-Noetherian-permanence}. If $I=0$, the empty generator
list gives $S=R$ in degree zero and the same argument applies.

The graded $S$-module
$E=\bigoplus_{n\geq0}I^nM$ has action
$\iota_d(a)\jmath_n(x)=\jmath_{d+n}(ax)$, where $\jmath_n$
denotes its degree inclusion. If $x_1,\ldots,x_r$ generate
$M$, their degree-zero images generate $E$: each member of
$I^nM$ is a sum $\sum_{i=1}^r a_ix_i$ with $a_i\in I^n$.
The graded subgroup
$$
L=\bigoplus_{n\geq0}(I^nM\cap N)\subset E
$$
is an $S$-submodule, since
$I^d(I^nM\cap N)\subset I^{d+n}M\cap N$.
By Lemma \ref{lemma-Noetherian-basic}, $L$ is finite over $S$.
Take a finite generating list and replace it by the collection of
all its homogeneous components. Each component is in $L$, each
old generator is the sum of its components, and the generated
submodule is consequently still exactly $L$. Denote this finite
list by $\jmath_{d_j}(\xi_j)$ for $1\leq j\leq s$, where
$\xi_j\in I^{d_j}M\cap N$. The list may be empty when $L=0$.
Set
$$
c=\max\bigl(\{1\}\cup\{d_j:1\leq j\leq s\}\bigr)>0.
$$
For $n\geq c$, taking the degree-$n$ component of an expression
in these generators shows that every $x\in I^nM\cap N$ is
$$
x=\sum_{j=1}^s h_j\xi_j,\qquad h_j\in I^{n-d_j}.
$$
The degree inclusions above specify the graded expression from
which this equality in the original $M$ is obtained.
For each $j$, the ideal identity
$I^{n-d_j}=I^{n-c}I^{c-d_j}$ means that one can write
$h_j=\sum_{\ell=1}^{q_j}a_{j\ell}b_{j\ell}$ with
$a_{j\ell}\in I^{n-c}$ and $b_{j\ell}\in I^{c-d_j}$.
Each $b_{j\ell}\xi_j$ lies in $I^cM\cap N$, so the full
equality
$x=\sum_{j=1}^s\sum_{\ell=1}^{q_j}
a_{j\ell}(b_{j\ell}\xi_j)$ proves
$I^nM\cap N\subset I^{n-c}(I^cM\cap N)$.
The opposite inclusion holds because a product with a member of
$I^cM\cap N$ remains in $N$ and lies in $I^nM$.
If the generating list is empty then $L=0$, both modules in the
desired equality are zero, and the same choice $c=1$ is valid.

Every larger integer $d\geq c$ is also valid. Indeed, for $n\geq d$,
the equalities already proved at $n$ and at $d$ give
$$
I^nM\cap N=I^{n-c}(I^cM\cap N)
=I^{n-d}\bigl(I^{d-c}(I^cM\cap N)\bigr)
=I^{n-d}(I^dM\cap N).
$$
In particular, the two original filtrations on $N$ satisfy
$$
I^nN\subset I^nM\cap N\quad(n\geq0),\qquad
I^{k+c}M\cap N\subset I^kN\quad(k\geq0).
$$
The first containment follows from $N\subset M$; the second
uses the equality at $n=k+c$ and $I^cM\cap N\subset N$.
Thus the subspace $I$-adic topology from $M$ and the $I$-adic
topology on $N$ have cofinal neighbourhood systems, with the
displayed explicit index bound.
\end{proof}